%Version 3.1 December 2024
% See section 11 of the User Manual for version history
%
%%%%%%%%%%%%%%%%%%%%%%%%%%%%%%%%%%%%%%%%%%%%%%%%%%%%%%%%%%%%%%%%%%%%%%
%%                                                                 %%
%% Please do not use \input{...} to include other tex files.       %%
%% Submit your LaTeX manuscript as one .tex document.              %%
%%                                                                 %%
%% All additional figures and files should be attached             %%
%% separately and not embedded in the \TeX\ document itself.       %%
%%                                                                 %%
%%%%%%%%%%%%%%%%%%%%%%%%%%%%%%%%%%%%%%%%%%%%%%%%%%%%%%%%%%%%%%%%%%%%%

%%\documentclass[referee,sn-basic]{sn-jnl}% referee option is meant for double line spacing

%%=======================================================%%
%% to print line numbers in the margin use lineno option %%
%%=======================================================%%

%%\documentclass[lineno,pdflatex,sn-basic]{sn-jnl}% Basic Springer Nature Reference Style/Chemistry Reference Style

%%=========================================================================================%%
%% the documentclass is set to pdflatex as default. You can delete it if not appropriate.  %%
%%=========================================================================================%%

%%\documentclass[sn-basic]{sn-jnl}% Basic Springer Nature Reference Style/Chemistry Reference Style

%%Note: the following reference styles support Namedate and Numbered referencing. By default the style follows the most common style. To switch between the options you can add or remove “Numbered” in the optional parenthesis. 
%%The option is available for: sn-basic.bst, sn-chicago.bst%  
 
%%\documentclass[pdflatex,sn-nature]{sn-jnl}% Style for submissions to Nature Portfolio journals
%%\documentclass[pdflatex,sn-basic]{sn-jnl}% Basic Springer Nature Reference Style/Chemistry Reference Style
\documentclass[pdflatex,sn-mathphys-num, oneside]{sn-jnl}% Math and Physical Sciences Numbered Reference Style
%%\documentclass[pdflatex,sn-mathphys-ay]{sn-jnl}% Math and Physical Sciences Author Year Reference Style
%%\documentclass[pdflatex,sn-aps]{sn-jnl}% American Physical Society (APS) Reference Style
%%\documentclass[pdflatex,sn-vancouver-num]{sn-jnl}% Vancouver Numbered Reference Style
%%\documentclass[pdflatex,sn-vancouver-ay]{sn-jnl}% Vancouver Author Year Reference Style
%%\documentclass[pdflatex,sn-apa]{sn-jnl}% APA Reference Style
%%\documentclass[pdflatex,sn-chicago]{sn-jnl}% Chicago-based Humanities Reference Style

%%%% Standard Packages
%%<additional latex packages if required can be included here>

\usepackage{graphicx}%
\usepackage{multirow}%
\usepackage{amsmath,amssymb,amsfonts}%
\usepackage{amsthm}%
\usepackage{mathrsfs}%
\usepackage[title]{appendix}%
\usepackage{xcolor}%
\usepackage{textcomp}%
\usepackage{manyfoot}%
\usepackage{booktabs}%
\usepackage{algorithm}%
\usepackage{algorithmicx}%
\usepackage{algpseudocode}%
\usepackage{listings}%
\usepackage{geometry}
\geometry{
  reset,          
  a4paper,
  textwidth=190mm,% Forces the text width (A4 is 210mm wide - 10L - 10R = 190mm)
  textheight=267mm, % Forces the text height (A4 is 297mm tall - 15T - 15B = 267mm)
  left=10mm,
  right=10mm,
  top=15mm,
  bottom=15mm
}
%%%%

%%%%%=============================================================================%%%%
%%%%  Remarks: This template is provided to aid authors with the preparation
%%%%  of original research articles intended for submission to journals published 
%%%%  by Springer Nature. The guidance has been prepared in partnership with 
%%%%  production teams to conform to Springer Nature technical requirements. 
%%%%  Editorial and presentation requirements differ among journal portfolios and 
%%%%  research disciplines. You may find sections in this template are irrelevant 
%%%%  to your work and are empowered to omit any such section if allowed by the 
%%%%  journal you intend to submit to. The submission guidelines and policies 
%%%%  of the journal take precedence. A detailed User Manual is available in the 
%%%%  template package for technical guidance.
%%%%%=============================================================================%%%%

%% as per the requirement new theorem styles can be included as shown below
\theoremstyle{thmstyleone}%
\newtheorem{theorem}{Theorem}%  meant for continuous numbers
%%\newtheorem{theorem}{Theorem}[section]% meant for sectionwise numbers
%% optional argument [theorem] produces theorem numbering sequence instead of independent numbers for Proposition
\newtheorem{proposition}[theorem]{Proposition}% 
%%\newtheorem{proposition}{Proposition}% to get separate numbers for theorem and proposition etc.

\theoremstyle{thmstyletwo}%
\newtheorem{example}{Example}%
\newtheorem{remark}{Remark}%

\theoremstyle{thmstylethree}%
\newtheorem{definition}{Definition}%

\raggedbottom
%%\unnumbered% uncomment this for unnumbered level heads

% --- NEW SPACE-SAVING COMMANDS ---
\usepackage{microtype}      % 1. Optimizes horizontal text packing (Saves ~5% space)
\linespread{0.96}           % 2. Reduces vertical space between lines (Saves ~5% space)
\usepackage{newtxtext}      % 3. Ensures you are using the most modern, compact Times font
\usepackage{newtxmath}
% ---------------------------------

\begin{document}

\title[Article Title]{Are meteorological, agricultural, and satellite features from the CY-Bench Dataset good predictors of European Wheat Futures Price Changes?}

%%=============================================================%%
%% GivenName	-> \fnm{Joergen W.}
%% Particle	-> \spfx{van der} -> surname prefix
%% FamilyName	-> \sur{Ploeg}
%% Suffix	-> \sfx{IV}
%% \author*[1,2]{\fnm{Joergen W.} \spfx{van der} \sur{Ploeg} 
%%  \sfx{IV}}\email{iauthor@gmail.com}
%%=============================================================%%

\author{\fnm{68729}}
\author{\fnm{candidate number}}
\author{\fnm{candidate number}}

%%\pacs[JEL Classification]{D8, H51}

%%\pacs[MSC Classification]{35A01, 65L10, 65L12, 65L20, 65L70}

\maketitle

\clearpage

\section{Executive Summary}\label{sec1}

\clearpage

\section{Table of Contents}\label{sec2}

\clearpage

\section{Introduction (1 page)}\label{sec3}

\clearpage

\section{(Optional) Literature Review (1 page)}\label{sec3}

\clearpage

\section{Data Pre-processing (2 pages)}\label{sec3}

\clearpage

\section{Redefining the Prediction Task (1 page)}\label{sec3}

\clearpage

\section{Feature Engineering and Selection (2 pages)}\label{sec3}

\clearpage

\section{Method 1 (2 pages)}\label{sec3}


\clearpage

\section{Method 2 (2 pages)}\label{sec3}
Random Forest and ExtraTrees are evaluated together as the tree-based ensemble methods in our pipeline, following directly from the linear baseline established by Ridge and Lasso. Both methods are motivated by the same core limitation identified in the previous sections finding that the linear models cannot represent interactions between features that together signal crop stress in ways that no single feature captures alone. Tree-based ensembles address this by partitioning the feature space through successive splits, naturally capturing these non-linear interactions without requiring them to be specified manually.

For the Futures Price Prediction Task, Random Forest (RF) is chosen for three reasons. With its bagging approach for training each tree on a bootstrapped sample of observations and averaging predictions across trees, we can aim to reduce variance relative to a single tree, which is particularly valuable in a noisy financial return prediction task where individual trees would overfit heavily to the training window. The random feature subsampling at each split de-correlates the trees, preventing any single dominant agricultural feature from driving all predictions and encouraging the ensemble to exploit the full breadth of the 76 available features. From a practical perspective, RF's robustness to hyperparameter choice makes it a reliable step up from linear models without requiring extensive tuning to obtain useful results. Aditionally, RF appears prominently in the yield prediction literature we reviewed~\cite{bib14,bib16,bib17}, making it a natural methodological bridge between the crop yield benchmarks and our futures application.

ExtraTrees is evaluated alongside RF as a direct comparison. The key distinction is that ExtraTrees uses random rather than optimal split thresholds at each node, introducing additional randomisation that further reduces variance at the cost of some increase in bias. In a setting with many weakly informative features and a noisy target, this extra randomisation may be beneficial by preventing the ensemble from overcommitting to any particular feature split that happens to perform well in training but does not generalise during testing. 

The hyperparameter search for both models covers n\_estimators from [200,500] for RF and fixed at 500 for ExtraTrees. max\_depth is set to set of 3,5, or 8, min\_samples\_leaf at 10, 20, or 40, with max\_features set to square root, log base 2, 0.3, or 0.5. All values are discrete rather than continuous because ensemble behaviour is not sensitive to small perturbations in these parameters. The meaningful distinctions are between shallow and deep trees, few and many estimators, and narrow versus wide feature subsets. The max\_depth upper bound of 8 prevents deep trees that would memorise individual training windows, which is critical in a walk-forward scheme where the training set grows by only one year at each step. The min\_samples\_leaf lower bound of 10 enforces a minimum leaf size that prevents extremely specific splits on small subsets of the 76 agricultural features. Both models use 20 random draws over the validation period from 2017 through 2019 with directional accuracy as the selection criterion for the same reasons we used directional accuracy as stated in the Ridge and Lasso Regression section.

The optimal RF configuration is n\_estimators=500, max\_depth=5, min\_samples\_leaf=40, max\_features=0.5, achieving a validation directional accuracy of 49.28\%. The large min\_samples\_leaf of 40 and shallow max\_depth of 5 together indicate that strong regularisation is needed even for tree ensembles on this task. This shows the model benefits from broad, large splits rather than finer grained partitions, consistent with the weak nature of the agricultural signal. For ExtraTrees, the optimal configuration is n\_estimators=500, max\_depth=5, min\_samples\_leaf=20, max\_features=sqrt, achieving a validation directional accuracy of 46.80\%. The lower min\_samples\_leaf relative to RF is consistent with ExtraTrees additional randomisation at split points providing implicit regularisation, reducing the need to enforce large leaf sizes explicitly.

Both models are fit using the same walk-forward expanding window scheme as the linear models, trained independently across all 515 regions with a Standard Scaler on training data only and predictions aggregated via production-weighted averaging.

On the test set from 2020 through 2023, RF achieves R\textsuperscript{2} = −0.0099, directional accuracy of 48.43\%, and RMSE of 0.0407. ExtraTrees achieves R\textsuperscript{2} = −0.0074, directional accuracy of 50.98\%, and RMSE of 0.0407. Both models improve over Ridge and Lasso on every metric. Notably, ExtraTrees is the only model in our entire pipeline to exceed 50\% directional accuracy on the test set, suggesting that the additional randomisation in ExtraTrees split selection provides a marginal but real advantage in this noisy prediction environment. However, both R\textsuperscript{2} values remain negative, confirming that neither ensemble model outperforms the unconditional mean on the futures task, similarly to Ridge and Lasso. The RMSE values are the same when rounded to two decimal places, and nearly the same to those of the linear models, again reflecting the dominant effect of the high-volatility 2022 period rather than model-specific behaviour.

For the Yield Task, RF and ExtraTrees are tuned separately, per country, over the validation period from 2013 to 2015 with region NRMSE as the scoring criterion, using the same hyperparameter grid as for the futures task but with max\_depth extended to include 12 and None (unlimited) to allow the trees to grow deeper on the tabular annual yield dataset, which has more structured signal than daily futures returns.

The optimal configurations differs meaningfully across countries. For Germany, RF selects n\_estimators=500, max\_depth=8, min\_samples\_leaf=20, max\_features=0.5, achieving a validation NRMSE of 11.36\%. ExtraTrees selects n\_estimators=500, max\_depth=12, min\_samples\_leaf=5, max\_features=0.5 with validation NRMSE of 11.76\%. Both models choose deeper trees and smaller leaf sizes compared to the futures task, reflecting the more structured annual yield signal where more precise feature interactions improve prediction. For France and Poland, both models independently select max\_depth=None and min\_samples\_leaf=2, allowing unrestricted tree growth. This may reflect that France's 90 regions and Poland's 16 regions provide smaller training sets where deeper trees do not overfit as severely, or that the yield signal in these countries is concentrated in strong feature interactions that require deeper splits to capture. RF acheives a regional NRMSE of 12.25\% and ExtraTrees is at 12.66\% on France. For Poland, RF is 17.43\% and ExtraTrees is 17.40\%. The higher validation NRMSE for France and Poland compared to Germany, combined with the need for unrestricted tree depth, suggests that the ensemble models are less effective in settings with fewer regions and less structured training data.

Both models are evaluated on the test period of 2016 through 2020 in a walk-forward scheme, training on all prior region-years with predictions aggregated to national level using production weights.

At end-of-season, both tree ensembles substantially outperform the linear models for Germany. RF achieves a region NRMSE of 13.57\% and EU NRMSE of 4.57\%, while ExtraTrees achieves 13.00\% and 4.24\% respectively, making ExtraTrees the best-performing model for Germany across all methods. These results are broadly comparable to Paudel et al.~\cite{bib16}'s German baseline of approximately 15\% regional NRMSE and confirm that non-linear interactions between weather, soil, and vegetation features are important for German yield prediction. For France, RF achieves region NRMSE of 17.97\% and EU NRMSE of 14.34\%, while ExtraTrees achieves 17.77\% and 13.90\%, again with ExtraTrees narrowly ahead. Both remain below the Paudel et al.~\cite{bib16} benchmark of approximately 6\% national NRMSE for the same WOFOST feature quality reasons discussed in the previous section. For Poland, RF achieves region NRMSE of 16.90\% and EU NRMSE of 12.34\%, while ExtraTrees achieves 16.24\% and 12.12\%. Both are substantially worse than Ridge's 12.61\% regional and 4.08\% EU result, confirming that Ridge's Poland performance was driven by the small sample size rather than a genuine linear signal advantage.

Across lead times, both tree models show the same gradual degradation from end-of-season toward earlier lead times that we observed for the linear models, consistent with Paudel et al.~\cite{bib16}'s finding that structural yield drivers dominate throughout the season. For Germany, ExtraTrees moves from 13.00\% at end-of-season to 14.60\% at mid-season and 15.53\% at quarter-season, a degradation of only 2.5 percentage points across all three lead times. RF follows an almost identical pattern at 13.57\%, 14.56\%, and 15.66\%. For France, ExtraTrees moves from 17.77\% to 19.12\% to 20.23\%, and RF from 17.97\% to 20.29\% to 21.26\%. Again, a gradual and consistent increase. Poland shows sharper degradation for both tree models, similar to the pattern seen with linear models and reinforcing the small-sample instability issue. Across all countries and lead times, ExtraTrees outperforms or matches RF, confirming that additional split randomisation is consistently beneficial for this yield prediction task, as illustrated in Figure~\ref{fig:NRMSE_results}. The year-by-year national aggregate predictions are shown in Figure~\ref{fig:yield_forecast_accuracy_graph}.

Overall, the tree-based ensembles establish a substantially improved baseline over the linear models for yield prediction, while delivering only marginal improvements on the futures task. The gap between the two tasks reinforces the core finding that agricultural features carry predictive information about crop yields but that information does not translate into exploitable futures price signal, consistent with semi-strong market efficiency. The neural network is examined next to assess whether greater model expressiveness can close this gap further.



\clearpage

\section{Method 3 (2 pages)}\label{sec3}

\clearpage

\section{Method 4 (2 pages)}\label{sec3}

\clearpage

\section{Model Result Comparison (1 page)}\label{sec3}

\clearpage

\section{Conclusion (1 page)}\label{sec3}

\clearpage

\bibliography{sn-bibliography}% common bib file
%% if required, the content of .bbl file can be included here once bbl is generated
%%\input sn-article.bbl

\clearpage

\section{Generative AI Statement}

\clearpage

\begin{appendices}

\section{Section title of first appendix}\label{secA1}

An appendix contains supplementary information that is not an essential part of the text itself but which may be helpful in providing a more comprehensive understanding of the research problem or it is information that is too cumbersome to be included in the body of the paper.

%%=============================================%%
%% For submissions to Nature Portfolio Journals %%
%% please use the heading ``Extended Data''.   %%
%%=============================================%%

%%=============================================================%%
%% Sample for another appendix section			       %%
%%=============================================================%%

%% \section{Example of another appendix section}\label{secA2}%
%% Appendices may be used for helpful, supporting or essential material that would otherwise 
%% clutter, break up or be distracting to the text. Appendices can consist of sections, figures, 
%% tables and equations etc.

\end{appendices}

%%===========================================================================================%%
%% If you are submitting to one of the Nature Portfolio journals, using the eJP submission   %%
%% system, please include the references within the manuscript file itself. You may do this  %%
%% by copying the reference list from your .bbl file, paste it into the main manuscript .tex %%
%% file, and delete the associated \verb+\bibliography+ commands.                            %%
%%===========================================================================================%%

\end{document}
